\subsection{Hodnotenie udalostí}                                                                                                                                                                                                 
                                                                                                                                                                                                                                   
  Systém hodnotenia umožňuje účastníkom ohodnotiť udalosť po jej skončení.                                                                                                                                                         
  Hodnotenie je implementované priamo v \texttt{AppViewModel}.                                                                                                                                                                                                                                                                                                                                                                                                        
  \subsubsection{Dátový model:}                                                                                                                                                                                                    
  Hodnotenie je uložené priamo v dokumente eventu v troch poliach:                                                                                                                                                                                                                                                                                                                                                                                                    
  \begin{lstlisting}[language=Kotlin, caption=Polia hodnotenia v dátovom modeli eventu]                                                                                                                                            
  val rating      : Double       = 0.0         // vážený priemer
  val totalRatings: Int          = 0           // počet hodnotení
  val ratedBy     : List<String> = emptyList() // UIDs ktorí už hodnotili
  \end{lstlisting}

  \subsubsection{Logika hodnotenia:}
  Nový priemer sa vypočíta ako vážený priemer — starý priemer sa vynásobí počtom
  predchádzajúcich hodnotení, pripočíta sa nové hodnotenie a výsledok sa vydelí
  novým celkovým počtom. Týmto spôsobom sa zachová presnosť bez nutnosti ukladať



  \subsubsection{Pravidlá hodnotenia:}
  \begin{itemize}[noitemsep, nolistsep]
      \item Hodnotiť môžu len účastníci eventu - nie tvorca ani cudzí používatelia.
      \item Každý používateľ môže hodnotiť len raz - kontrola prebieha cez zoznam
            \texttt{ratedBy}.
      \item Hodnotenie je v rozsahu 1 až 5 hviezd.
      \item Priemerné hodnotenie vstupuje ako jeden z faktorov do odporúčacieho
            systému s váhou 16 ak ju používateľ nezmení\,\%.
  \end{itemize}

  \subsubsection{Používateľské rozhranie:}
  Sekcia hodnotenia sa v \texttt{EventDetailScreen} zobrazí len
  účastníkovi, ktorý ešte nehodnotil. Aktuálne priemerné hodnotenie a celkový
  počet hodnotení sú viditeľné pre všetkých používateľov.

  \iffalse                                                                                                                                                                                                                         
  \begin{lstlisting}[language=Kotlin, caption=Podmienka zobrazenia sekcie hodnotenia]
  val isAttending = currentUserId in event.attendees
  val hasRated    = currentUserId in event.ratedBy

  if (isAttending && !isCreator && !hasRated) {
  // 5 hviezd - Row s ikonami
  }
  \end{lstlisting}
  \fi                                                                                                                                                                                                                              

  % -------------------------------------------------------                                                                                                                                                                        
  \subsubsection{Check-in systém}

  Check-in systém umožňuje organizátorovi overiť fyzickú prítomnosť účastníkov
  na udalosti. Pri tvorbe eventu si organizátor zvolí jeden zo štyroch typov
  check-inu, ktoré sú zhrnuté v tabuľke~\ref{tab:checkin-types}.

  \begin{table}[h]
      \centering                                                                                                                                                                                                                   
      \caption{Typy check-inu}
      \label{tab:checkin-types}
      \begin{tabular}{ll}
          \toprule                                                                                                                                                                                                                 
          \textbf{Hodnota} & \textbf{Popis} \\                                                                                                                                                                                     
          \midrule                                                                                                                                                                                                                 
          \texttt{none}            & Check-in sa nevyžaduje \\                                                                                                                                                                     
          \texttt{qr}              & Účastník naskenuje QR kód \\                                                                                                                                                                  
          \texttt{manual}          & Organizátor manuálne označí prítomnosť \\                                                                                                                                                     
          \texttt{qr\_and\_manual} & Kombinácia oboch predchádzajúcich \\                                                                                                                                                          
          \bottomrule                                                                                                                                                                                                              
      \end{tabular}
  \end{table}

  \iffalse                                                                                                                                                                                                                         
  \paragraph{Generovanie QR kódu:}
  QR kód sa generuje pomocou knižnice ZXing priamo na zariadení organizátora
  (\texttt{QrCodeScreen.kt}). Obsahom QR kódu je \texttt{eventId} — unikátny
  identifikátor udalosti z Firestore:

  \begin{lstlisting}[language=Kotlin, caption=Generovanie QR kódu pomocou ZXing]
  private fun generateQrBitmap(content: String, size: Int = 512): Bitmap? {
      val bitMatrix = QRCodeWriter()
          .encode(content, BarcodeFormat.QR_CODE, size, size)
      val bitmap = Bitmap.createBitmap(size, size, Bitmap.Config.RGB_565)
      for (x in 0 until size) {
          for (y in 0 until size) {
              bitmap.setPixel(x, y,
                  if (bitMatrix[x, y]) Color.BLACK else Color.WHITE
              )
          }
      }
      return bitmap
  }
  \end{lstlisting}
  \fi                                                                                                                                                                                                                              

  \paragraph{Generovanie QR kódu:}
  QR kód sa generuje pomocou knižnice ZXing priamo na zariadení organizátora
  (\texttt{QrCodeScreen.kt}). Obsahom kódu je \texttt{eventId} - unikátny
  identifikátor udalosti z Firestore.

  \paragraph{Skenovanie QR kódu:}
  Skenovanie je implementované v \texttt{QrScanScreen.kt} pomocou CameraX
  na zobrazenie náhľadu kamery a ML Kit Barcode Scanning na rozpoznanie kódu.
  Obrazovka najprv požiada o oprávnenie . Stratégia
  \texttt{STRATEGY\_KEEP\_ONLY\_LATEST} zabezpečuje spracovanie len posledného
  snímku. Po úspešnom naskenovaní sa zavolá
  \texttt{appViewModel.checkInAttendee(eventId, userId)}.

  \iffalse                                                                                                                                                                                                                         
  \begin{lstlisting}[language=Kotlin, caption=Skenovanie QR kódu pomocou ML Kit]
  val barcodeScanner = BarcodeScanning.getClient(
      BarcodeScannerOptions.Builder()
          .setBarcodeFormats(Barcode.FORMAT_QR_CODE)
          .build()
  )

  val analysis = ImageAnalysis.Builder()
      .setBackpressureStrategy(ImageAnalysis.STRATEGY_KEEP_ONLY_LATEST)
      .build()

  analysis.setAnalyzer(ContextCompat.getMainExecutor(ctx)) { imageProxy ->
      val image = InputImage.fromMediaImage(
          mediaImage, imageProxy.imageInfo.rotationDegrees
      )
      barcodeScanner.process(image)
          .addOnSuccessListener { barcodes ->
              barcodes.firstOrNull()?.rawValue?.let { value ->
                  scanned = true
                  onScanResult(value)  // value = eventId
              }
          }
  }
  \end{lstlisting}
  \fi                                                                                                                                                                                                                              

  \paragraph{Check-in logika:}
  Metóda \texttt{checkInAttendee()} overí, či je používateľ účastníkom eventu
  a či ešte nebol check-inovaný:

  \begin{lstlisting}[language=Kotlin, caption=Logika check-inu v AppViewModel]
  fun checkInAttendee(eventId: String, userId: String) {
      val event = _events[eventId] ?: return
      if (userId !in event.attendees) {
          checkInError = "Nie si účastníkom tejto udalosti."
          return
      }
      if (userId in event.checkedIn) {
          checkInError = "Už si zaregistrovaný."
          return
      }
      val updated = event.copy(checkedIn = event.checkedIn + userId)
      _events[eventId] = updated
      viewModelScope.launch {
          db.collection("events").document(eventId)
              .update("checkedIn", updated.checkedIn).await()
          checkInSuccess = true
      }
  }
  \end{lstlisting}

  Organizátor môže check-in aj zrušiť volaním \texttt{checkOutAttendee()},
  ktorá odstráni \texttt{userId} zo zoznamu \texttt{checkedIn}.

  \paragraph{Manuálny check-in:}
  Pri type \texttt{manual} alebo \texttt{qr\_and\_manual} vidí organizátor
  v Tab~Účastníci pri každom účastníkovi prepínač $\circ$\,/\,$\checkmark$.
  Kliknutím sa zavolá \texttt{checkInAttendee()} alebo \texttt{checkOutAttendee()}.
  Účastník s potvrdenou prítomnosťou má zelenú odznak „$\checkmark$ Prítomný".

  \paragraph{Väzba na penalizačný systém:}
  Ak je \texttt{checkInType} \texttt{none} a účastník sa neodhlási ani
  nezaregistruje check-in, po skončení eventu mu je automaticky zaznamenaný
  no-show.

  % -------------------------------------------------------                                                                                                                                                                        
  \subsubsection{Systém penalizácií (No-show)}

  Penalizačný systém motivuje používateľov dodržiavať záväzky voči udalostiam.
  Ak sa používateľ prihlási na udalosť s povoleným check-inom, fyzicky sa
  nezúčastní a ani sa včas neodhlási, systém mu zaznamená no-show.

  \paragraph{Dátový model:}
  Penalizácie sú uložené priamo v profile používateľa (\texttt{UserProfile}):

  \begin{lstlisting}[language=Kotlin, caption=Polia no-show a soft ban v UserProfile]
  val noShowTimestamps: List<Long>   = emptyList() // čas každého no-show (epoch ms)
  val noShowEventIds  : List<String> = emptyList() // dedup - eventy kde bol no-show

  val banned          : Boolean      = false        // true = účet zablokovaný
  val banReason       : String       = ""           // "no_show" | "reports" | ""
  val banTimestamp    : Long         = 0L           // čas banu v epoch ms
  \end{lstlisting}

  Systém implementuje \textit{soft ban} --- pri zablokovaní sa používateľské dáta
  (história, priatelia, správy) nemažú. Administrátor môže priamo vo Firebase
  Console vidieť dôvod (\texttt{banReason}) a čas (\texttt{banTimestamp}) banu
  a manuálne rozhodnúť o odbanovaní alebo permanentnom vymazaní účtu.

  \iffalse                                                                                                                                                                                                                         
  Pomocná metóda počíta no-show len za posledných 30 dní:

  \begin{lstlisting}[language=Kotlin, caption=Výpočet počtu no-show za posledných 30 dní]
  fun recentNoShowCount(): Int {
      val cutoff = System.currentTimeMillis() - 30L * 24 * 3_600_000
      return noShowTimestamps.count { it > cutoff }
  }
  \end{lstlisting}
  \fi                                                                                                                                                                                                                              

  \paragraph{Dva spôsoby vzniku no-show:}
  \begin{enumerate}
      \iffalse                                                                                                                                                                                                                     
      \item \textbf{Neskoré odhlásenie} — používateľ sa odhlási po uplynutí
            \texttt{cancelDeadlineHours}:
  \begin{lstlisting}[language=Kotlin]
  fun leaveEvent(eventId: String, userId: String) {
      val applyPenalty = event.isPastCancelDeadline()
      // ... odhlásenie z eventu ...
      if (applyPenalty) recordNoShow(userId, eventId)
  }
  \end{lstlisting}
      \fi                                                                                                                                                                                                                          
      \item \textbf{Neskoré odhlásenie} — používateľ sa odhlási po uplynutí
            \texttt{cancelDeadlineHours} (podrobnejšie popísané v sekcii
            správy účasti).

      \item \textbf{Neúčasť} — systém pri štarte skontroluje všetky minulé
            eventy, kde bol používateľ prihlásený, ale nebol check-inovaný:
  
  \end{enumerate}


  \paragraph{Zablokovanie účtu:}
  Ak ná používateľ 3 zablokovania dostáva ban. Pri ďalšom prihlásení \texttt{loadUserProfile()} toto
  pole skontroluje a okamžite odhlási používateľa — rovnaký mechanizmus je
  popísaný v sekcii autentifikácie.

  \iffalse                                                                                                                                                                                                                         
  \begin{lstlisting}[language=Kotlin, caption=Kontrola zablokovania pri prihlásení]
  if (profile.banned) {
      FirebaseAuth.getInstance().signOut()
      loginBannedError = true
      return@launch
  }
  \end{lstlisting}
  \fi                                                                                                                                                                                                                              

  No-show sa počítajú výlučne za posledných 30 dní — po uplynutí tohto obdobia
  sa staré záznamy nezapočítavajú, čo umožňuje rehabilitáciu používateľa bez
  nutnosti manuálneho zásahu administrátora.

  % -------------------------------------------------------                                                                                                                                                                        
  \subsubsection{Systém nahlasovania}

  Systém nahlasovania umožňuje používateľom upozorniť na nevhodný obsah -
  či už udalosť alebo iného používateľa. Pri dosiahnutí stanoveného počtu
  nahlásení sa obsah automaticky odstráni bez nutnosti zásahu administrátora.

  \iffalse                                                                                                                                                                                                                         
  \paragraph{Dátový model:}
  Každé nahlásenie je reprezentované triedou \texttt{Report} a uložené
  do Firestore kolekcie \texttt{reports}:

  \begin{lstlisting}[language=Kotlin, caption=Dátová trieda Report]
  data class Report(
      val id        : String = "",
      val reporterId: String = "",
      val targetId  : String = "",
      val targetType: String = "",
      val reason    : String = "",
      val timestamp : Long   = System.currentTimeMillis()
  )
  \end{lstlisting}
  \fi                                                                                                                                                                                                                              

  Každé nahlásenie je uložené do Firestore kolekcie \texttt{reports} s poliami
  \texttt{reporterId}, \texttt{targetId}, \texttt{targetType} (\texttt{"event"}
  alebo \texttt{"user"}), \texttt{reason} a \texttt{timestamp}.

  \paragraph{Dôvody nahlasovania:}
  Zoznam dôvodov a hranica automatickej akcie sú definované ako konštanty
  v \texttt{AppViewModel}:

  \begin{lstlisting}[language=Kotlin, caption=Konštanty systému nahlasovania]
  companion object {
      val REPORT_REASONS = listOf(
          "Nevhodný obsah",
          "Spam",
          "Falošná udalosť / profil",
          "Obťažovanie",
          "Iné"
      )
      private const val REPORTS_TO_BAN = 3
  }
  \end{lstlisting}

 
  


  \paragraph{Logika nahlasovania používateľa:}
  Metóda \texttt{reportUser()} funguje rovnakým spôsobom. Pri dosiahnutí
  prahu sa namiesto mazania udalosti nastaví príznak zablokovania
  na profile používateľa:

  \begin{lstlisting}[language=Kotlin, caption=Automatické zablokovanie používateľa]
  if (total.size() >= REPORTS_TO_BAN) {
      val banTime = System.currentTimeMillis()
      db.collection("users")
        .document(targetUserId)
        .update(mapOf(
            "banned"       to true,
            "banReason"    to "reports",
            "banTimestamp" to banTime
        )).await()
      cancelPendingInvitations(targetUserId)
  }
  \end{lstlisting}

  \iffalse                                                                                                                                                                                                                         
  \paragraph{Zrušenie čakajúcich pozvánok pri bane:}
  Pri každom zablokovaní účtu sa automaticky vymažú všetky čakajúce pozvánky
  adresované zabanovanému používateľovi:

  \begin{lstlisting}[language=Kotlin, caption=Implementácia cancelPendingInvitations]
  suspend fun cancelPendingInvitations(userId: String) {
      val pending = db.collection("invitations")
          .whereEqualTo("receiverId", userId)
          .whereEqualTo("status", "pending")
          .get().await()
      pending.documents.forEach { it.reference.delete().await() }
  }
  \end{lstlisting}
  \fi                                                                                                                                                                                                                              

  Pri každom zablokovaní (no-show aj nahlásenia) sa automaticky vymažú všetky
  čakajúce pozvánky adresované zabanovanému používateľovi volaním
  \texttt{cancelPendingInvitations()}.

  \paragraph{Pravidlá systému:}
  Tabuľka~\ref{tab:report-rules} sumarizuje pravidlá platné pre systém
  nahlasovania.

  \begin{table}[h]
      \centering                                                                                                                                                                                                                   
      \caption{Pravidlá systému nahlasovania}
      \label{tab:report-rules}
      \begin{tabular}{lp{8.5cm}}
          \toprule                                                                                                                                                                                                                 
          \textbf{Pravidlo} & \textbf{Detail} \\                                                                                                                                                                                   
          \midrule                                                                                                                                                                                                                 
          Kto môže nahlásiť &                                                                                                                                                                                                      
              Každý prihlásený používateľ okrem tvorcu udalosti
              a seba samého. \\                                                                                                                                                                                                    
          Duplikát &                                                                                                                                                                                                               
              Jeden používateľ môže nahlásiť ten istý cieľ len raz. \\                                                                                                                                                             
          Prah &                                                                                                                                                                                                                   
              3 unikátne nahlásenia spustia automatickú akciu. \\                                                                                                                                                                  
          Udalosť &                                                                                                                                                                                                                
              Vymazaná z Firestore aj z in-memory mapy. \\                                                                                                                                                                         
          Používateľ &                                                                                                                                                                                                             
              \texttt{banned = true} $\rightarrow$ odhlásenie
              pri ďalšom prihlásení. \\                                                                                                                                                                                            
          \bottomrule                                                                                                                                                                                                              
      \end{tabular}
  \end{table}

  % -------------------------------------------------------                                                                                                                                                                        
  \subsection{Implementácia komunikačného modulu v reálnom čase}

  Základom interaktivity je integračná vrstva služby Firebase Cloud
  Firestore, ktorá umožňuje obojsmernú synchronizáciu dát v reálnom čase
  bez nutnosti vlastnej serverovej infraštruktúry a správy WebSocket spojení.

  \iffalse                                                                                                                                                                                                                         
  \subsubsection{Dátový model správy}
  Pre potreby chatu bol definovaný dátový model \texttt{ChatMessage}:

  \begin{lstlisting}[language=Kotlin, caption={Dátový model správy v jazyku Kotlin}]
  data class ChatMessage(
      val id: String = "",
      val senderId: String = "",
      val senderName: String = "",
      val text: String = "",
      val timestamp: Long = 0L
  )
  \end{lstlisting}
  \fi                                                                                                                                                                                                                              

  Správy sú uložené v podkolekcii \texttt{events/\{eventId\}/messages/}
  s poliami \texttt{senderId}, \texttt{senderName}, \texttt{text}
  a \texttt{timestamp}.

  \subsubsection{Reaktívna synchronizácia dát}
  Technické jadro komunikácie tvorí metóda \texttt{addSnapshotListener}.
  Na rozdiel od štandardných HTTP požiadaviek vytvára trvalé spojenie
  s podkolekciou správ. Akonáhle dôjde k zápisu novej správy, Firebase
  iniciuje push notifikáciu a aplikácia aktualizuje reaktívny stav
  \texttt{chatMessages} vo ViewModeli.

  \begin{lstlisting}[language=Kotlin, caption={Implementácia real-time listenera vo ViewModeli}]
  fun startChatListener(eventId: String) {
      chatListener?.remove()
      chatListener = db.collection("events").document(eventId)
          .collection("messages")
          .orderBy("timestamp")
          .addSnapshotListener { snapshot, _ ->
              _chatMessages.clear()
              snapshot?.documents?.forEach { doc ->
                  doc.toObject(ChatMessage::class.java)?.let {
                      _chatMessages.add(it)
                  }
              }
          }
  }
  \end{lstlisting}